\chapter{Background}
\label{ch:background}

This chapter defines the terms used in the thesis and gives the legal and
conceptual background of the study.

\section{Key Terms}
\label{sec:terms}
\begin{description}[leftmargin=1.2em,labelindent=0pt,itemsep=2pt,
                    font=\normalfont\bfseries]
  \item[Health tracking app.] An app that lets people record data about their
    own health or body over time, for example steps and workouts, meals and
    weight, sleep, heart rate, mood or the menstrual cycle. The data come from
    the user's own entries, from the phone's sensors or from a paired device
    such as a fitness band. The thesis also calls such an app a
    \emph{health tracker}.
  \item[Health data.] The General Data Protection Regulation (GDPR) defines
    data concerning health as personal data about a person's physical or
    mental health that reveal information about their health status
    (Art.~4(15)) and places them among the special categories of personal data
    (Art.~9(1))~\parencite{gdpr2016}. Tracker data often reveal more than a
    single entry suggests: activity, sleep and heart-rate records can be used
    to infer health conditions, habits, locations and even personality
    traits~\parencite[\S4]{niksirat2024wearable}.
  \item[Surface.] A place where a provider tells users how it handles their
    data. This thesis looks at two surfaces: the privacy policy and the
    interface of the app, such as consent screens, settings and short notices.
  \item[Privacy harm.] A negative consequence that a privacy violation has for
    a person. Citron and Solove sort privacy harms into seven types: physical,
    economic, reputational, psychological, autonomy, discrimination and
    relationship harms~\parencite{citron2021harms}.
  \item[Pathway to harm.] A chain that starts with parts of a system, such as
    data disclosed to another party, continues with a capability that these
    parts give someone, such as tracking or inference, and ends in a harm. A
    \emph{harm tree} combines the pathways that lead to one
    harm~\parencite[\S4.1]{knabenhans2026harms}.
  \item[Protection claim.] A statement by a provider that a measure protects
    users' data, for example encryption, pseudonymisation or processing on the
    device. Such measures are often called privacy-enhancing technologies
    (PETs).
\end{description}

\section{What Providers Must Tell Users}
\label{sec:transparency}
The GDPR requires that information about the processing of personal data is
concise, transparent, intelligible and written in clear and plain language
(Art.~12). It must name, among other things, the purposes of the processing,
the recipients of the data and the rights of the user (Art.~13 and~14).
Consequences must be explained only for automated decisions with legal or
similarly significant effects (Art.~13(2)(f) and~22)~\parencite{gdpr2016}.
Guidance of the Article~29 Working Party goes further. It allows providers to
spread the information over layers, such as a privacy statement and short
notices inside the app, and it recommends that for complex or unexpected
processing providers also explain the most important consequences for the
people concerned~\parencite[paras.~10, 35, 39]{wp292018transparency}. Outside
automated decisions, the law therefore asks for a description of data
practices, and an explanation of what could follow from them is a
recommendation.

\section{From Data Practices to Harms}
\label{sec:practices-harms}
Knowing what happens to their data does not tell users what can happen to
them. A policy may say that data are shared with advertising partners without
saying who could then learn about a pregnancy or an illness. Harm trees make
this link visible. Knabenhans et al.\ use them to analyse the European Digital
Identity Framework and find that PETs reduce information flows, but that the
pathways to harm that come from the core function of the system
remain~\parencite{knabenhans2026harms}. A protection claim can therefore be
true and still leave a harm possible. The thesis builds on this idea and compares
what health tracker providers say about harms and protection with the harms
that their apps make possible.
